% §4.3 Closed-Loop Protocol — draft (phase 3)
\subsection{Closed-Loop Protocol and Logging}
\label{sec:closed-loop}

Each trial executes \texttt{capture} $\rightarrow$ \texttt{edge\_iqa} $\rightarrow$ \texttt{route} $\rightarrow$ action. Wall-clock segment timestamps and decisions are appended as one JSON object per line (JSONL) for supplementary reproducibility and auxiliary track (M6) trend analysis.

\paragraph{JSONL schema.}
Required fields: \texttt{trial\_id}, \texttt{image\_path} (relative), \texttt{q\_img}, \texttt{flags}, \texttt{retry\_count}, \texttt{route\_decision}, \texttt{t\_capture}, \texttt{t\_iqa}, \texttt{t\_route}, \texttt{latency\_ms}. Validation is automated via \texttt{scripts/validate\_jsonl.py}.

\paragraph{Execution.}
The driver \texttt{python -m langgraph\_router.run} alternates synthetic clear/blur frames when TCM val is unavailable; \texttt{scripts/paper1\_closed\_loop.sh} wraps 10 trials into \texttt{experiments/logs/run\_001.jsonl}. Main quantitative RTT (M1) remains on offline \texttt{run\_matrix.py}; JSONL supports protocol demonstration only.

\begin{table}[t]
\centering
\caption{Closed-loop JSONL fields (Paper~I).}
\begin{tabular}{ll}
\hline
Field & Description \\
\hline
\texttt{t\_capture}, \texttt{t\_iqa}, \texttt{t\_route} & ISO~8601 UTC segment times \\
\texttt{route\_decision} & \texttt{upload\_cloud} $|$ \texttt{resample\_edge} $|$ \texttt{fail\_safe} \\
\texttt{latency\_ms} & End-to-end trial latency \\
\hline
\end{tabular}
\end{table}
